\section{Extracting primes from the marked family}\label{sec:prime-reservoir}

We retain the notation and weighted family of
Section~\ref{sec:marked-family}.  In particular, a supported integer has
the unique form $d=4rQ+1$, where $r$ is a group integer with
$r\le x^{0.11}$ and $Q>x^{0.9}$ is prime, and it lies in the reduced
class $a$ modulo $M$.  By Lemma~\ref{lem:marked-mass} and $w(d)\le B_K$,
it is enough to prove that the prime weight is $\gg X_0/L$.
The two sieve inputs needed below are recorded explicitly.

\begin{hypothesis}[Block sieve]\label{input:block}
Let $z>1$, let $\mathcal A$ be a finite set with nonnegative weights
$\lambda_\nu$ for $\nu\in\mathcal A$, and let $\mathcal D$ be a set of
designated primes at most $z$.  For each $p\in\mathcal D$ let
$\mathcal A_p\subseteq\mathcal A$ be its bad condition.  For a squarefree
product $\ell$ of primes in $\mathcal D$, put
$\mathcal A_\ell=\bigcap_{p\mid\ell}\mathcal A_p$, with
$\mathcal A_1=\mathcal A$.  Suppose that, including for $\ell=1$,
\[
 \sum_{\nu\in\mathcal A_\ell}\lambda_\nu=Yg(\ell)+R(\ell),
 \qquad Y\ge0,
\]
where $g$ is multiplicative on these squarefree products and $g(1)=1$.
Extend $g(p)$ by zero at primes outside $\mathcal D$.  Suppose that for
fixed constants $\eta_0>0$ and $C_0<\infty$,
\[
 0\le g(p)\le1-\eta_0,
 \qquad
 \sum_{u<p\le u^2}g(p)\le C_0\quad(u>1).
\]
Write
\[
 S=\sum_{\nu\in\mathcal A\setminus\bigcup_{p\in\mathcal D}\mathcal A_p}
       \lambda_\nu.
\]
For every sufficiently large even integer $H$, the weight avoiding all
the bad conditions is
\begin{equation}\label{eq:block-sieve}
 \begin{split}
 S={}&Y\prod_{p\le z}(1-g(p))\bigl(1+O_{\eta_0,C_0}(e^{-H})\bigr)\\
 &+O_{\eta_0,C_0}\left(
    \sum_{\substack{\ell\le z^{4H+2}\\
                    \ell\ \mathrm{squarefree}\\
                    p\mid\ell\Rightarrow p\in\mathcal D}}
       |R(\ell)|\right).
 \end{split}
\end{equation}
The threshold for $H$ and the implied constants depend
only on $\eta_0,C_0$.  In particular, the constants are uniform when
$H$ grows, and the remainder sum has no additional coefficient or
multiplicity factor.  For $H=2$ there is the upper bound
\begin{equation}\label{eq:block-sieve-upper}
 S\ll_{\eta_0,C_0}
 Y\prod_{p\le z}(1-g(p))+
    \sum_{\substack{\ell\le z^{10}\\
                    \ell\ \mathrm{squarefree}\\
                    p\mid\ell\Rightarrow p\in\mathcal D}}
       |R(\ell)|.
\end{equation}
Both remainder sums include only squarefree products of the designated
primes, with $\ell=1$ included.  This is the form stated in
\cite[Lemma~11.1, Equations~(11.2)--(11.3)]{PrimitiveRoot2}.
\end{hypothesis}

\begin{hypothesis}[Rough-number density]\label{input:density}
For the density $D_\gamma(w)$ defined in Section~\ref{sec:marked-family},
there are absolute constants $b_d,c_d,C_d>0$ such that, for every fixed
$0<b<\min(b_d,1/2)$,
\begin{equation}\label{eq:density-integral-bound}
 \int_b^{1/2}D_t(1-t)\frac{\dd t}{t}=D_b(1)-1,
 \qquad
 D_b(1)\le\frac{e^{-\gamma_E}}{b}
                 \bigl(1+C_de^{-c_d/b}\bigr).
\end{equation}
Here $\gamma_E$ is Euler's constant, and the value of the integrand at
$t=1/2$ is immaterial.  These are the assertions of
\cite[Lemma~11.2, Equations~(11.6)--(11.7)]{PrimitiveRoot2}.
\end{hypothesis}

The identity in \eqref{eq:density-integral-bound} describes removal of
the least factor.  In the $j$-factor summand of $D_b(1)$, $j\ge2$, the
measure on $t_1+\cdots+t_j=1$, $t_i\ge b$, is
$\frac1{j!}(t_1\cdots t_j)^{-1}\dd t_1\cdots\dd t_{j-1}$.
It is invariant under permutation of the coordinates.  Its minimum
$t$ is unique outside a null set and is at most $1/2$.  Choosing this
coordinate in $j$ ways leaves $\dd t/t$ times the $(j-1)$-factor
summand of $D_t(1-t)$, since $j/j!=1/(j-1)!$.  The omitted one-factor
summand of $D_b(1)$ is $1$.

Measured in units of $\mathfrak S_M X_0/L$, the initial sieve below has
main coefficient $e^{-\gamma_E}/b$, while the bins for the least prime factor cost at
most $D_b(1)-1$ apart from the partition and approximation errors.
The density bound and the exponentially small loss in the initial sieve leave a
positive constant for small fixed $b$; a separate sieve controls the
remaining balanced semiprimes.

\subsection{The initial rough counts}

For now take any fixed $0<\epsilon<\kappa<0.01$ as in
Section~\ref{sec:marked-family}, and introduce a sufficiently small
fixed $0<b<0.01$.  The choices of these parameters and of $K$ will be
made in their required order at the end.  Every estimate below is for
$x$ tending to infinity after its displayed parameters have been fixed.
Write
\[
 \Pi_r(y)=\prod_{\substack{2<p\le y\\p\nmid Mr}}
                 \left(1-\frac1{p-1}\right).
\]
By \eqref{eq:missing-group-factors},
$\Pi_r(y)=V_M(y)(1+o(1))$ uniformly in $y$ and in all supported
$r\le x^{0.11}$.  A supported $d$ is odd and is coprime to $M$.
Consequently, for supported $d$, the condition $P^-(d)>z$ is equivalent
to excluding divisibility by the odd primes $p\le z$ with $p\nmid M$.

Fix a group integer $r\le x^\epsilon$.  Apply Input~\ref{input:block}
to the primes $Q$ and their weights in $A_r$, with the bad condition
$p\mid4rQ+1$ at each odd $p\le x^b$ not dividing $M$.  The intersection
counts are $A_rg_r(\ell)+R_r(\ell)$ by
\eqref{eq:distribution-data}.  Compact support of $\Psi$ makes this
a finite weighted family after removing zero weights.  The densities
on designated primes
obey
\[
 0\le g_r(p)\le\frac1{p-1}\le\frac12.
\]
Their sums on $(u,u^2]$ have an absolute bound, since
$1/(p-1)\le2/p$ for odd $p$ and the prime harmonic sums on these
intervals are bounded.  Thus the local constants of the block sieve
are absolute, independently of $M,r,K$.

Use the even depth
\[
 H_b=2\left\lfloor\frac1{40b}\right\rfloor.
\]
For sufficiently small $b$ it is an admissible depth in
Input~\ref{input:block}, and
\begin{equation}\label{eq:initial-sieve-level}
 H_b\ge\frac1{20b}-2,
 \qquad
 b(4H_b+2)\le\frac15+2b<0.22<\frac12-\frac\kappa2.
\end{equation}
In particular, every remainder modulus is at most
$x^{b(4H_b+2)}\le D_0$.  Sum the resulting lower bound with the
nonnegative marks $\mathcal W(r)$.  By
\eqref{eq:averaged-distribution} the sum of the remainders is
$O(xJ_0L^{-A})$ for any fixed $A$; taking $A>2$ makes it $o(X_0/L)$
by \eqref{eq:marked-mass-comparability}.  The tail with
$r>x^\epsilon$ is negligible to this accuracy by
\eqref{eq:marked-tail}.  The uniform formula for $\Pi_r(x^b)$ and
\eqref{eq:progression-product-asymptotic} give
\[
 \Pi_r(x^b)\sim V_M(x^b)
       \sim\frac{\mathfrak S_M e^{-\gamma_E}}{bL}.
\]
Finally, \eqref{eq:initial-sieve-level} gives
$e^{-H_b}\le e^2e^{-1/(20b)}$.  We have proved
\begin{equation}\label{eq:initial-rough-mass}
 \sum_{\substack{d\ge2\\P^-(d)>x^b}}w(d)
 \ge\frac{\mathfrak S_M X_0}{L}
 \left\{\frac{e^{-\gamma_E}}{b}
              \bigl(1-C_se^{-c_s/b}\bigr)+o(1)\right\},
\end{equation}
where $C_s,c_s>0$ are absolute; for example, $c_s=1/20$ is permitted
after enlarging $C_s$.  The implied threshold may depend on all fixed
parameters.

For fixed $b\le\gamma<\gamma'\le1/2-\kappa$, let
\[
 \mathcal N_{\gamma,\gamma'}
     =\{n\text{ prime}:x^\gamma<n\le x^{\gamma'}\}.
\]
To bound the $W$-rough term in the Type II replacement for $d=mn$, we use
the following count conditioned on $n\mid d$:
\begin{equation}\label{eq:conditioned-bin-count}
 \sum_{n\in\mathcal N_{\gamma,\gamma'}}
 \ \sum_{\substack{d\ge2\\n\mid d\\P^-(d)>W}}w(d)
 =X_0V_M(W)\left\{\log\frac{\gamma'}\gamma+o(1)\right\}.
\end{equation}
For large $x$, every such $n$ exceeds $W$, every prime dividing $M$,
and every group prime.  Indeed $x^b$ eventually exceeds each of these
quantities.  In particular, $n\nmid r$ for every group integer $r$.
At fixed $r\le x^\epsilon$, condition the family of $A_r$ on
$n\mid4rQ+1$ and designate the odd primes $p\le W$ not dividing $M$.
For a squarefree product $\ell$ of these primes, including $\ell=1$,
\begin{equation}\label{eq:conditioned-intersections}
 A_r(n\ell)=\frac{A_r}{n-1}g_r(\ell)+R_r(n\ell).
\end{equation}
This follows from $n\nmid Mr\ell$ and multiplicativity of $g_r$.
Thus the block sieve applies with $Y=A_r/(n-1)$, the same local
densities, and remainder $R_r(n\ell)$.  Its remainder at $\ell=1$ is
$R_r(n)$, which is explicitly allowed by Input~\ref{input:block}.

Take $H_W=2\lceil(\log L)^2\rceil$.  It is an admissible growing even
depth, its relative error tends to zero, and
\[
 W^{4H_W+2}
    =\exp\bigl(O(L^{0.24}(\log L)^2)\bigr)=x^{o(1)}.
\]
Every remainder modulus in \eqref{eq:conditioned-intersections}
therefore satisfies
\[
 n\ell\le x^{1/2-\kappa+o(1)}<x^{1/2-\kappa/2}=D_0
\]
for sufficiently large $x$.  These moduli are odd squarefree and
coprime to $M$.  Moreover the map $(n,\ell)\mapsto n\ell$ is
injective throughout the sum: $n$ is the unique prime factor of this
product exceeding $W$.  Hence the unweighted remainder sums of the
block sieve, summed over $n$ and then with the marks over $r$, form a
subsum of \eqref{eq:averaged-distribution}, without any multiplicity
or depth factor.  Taking its accuracy $A>2$ makes this
$o(X_0V_M(W))$, because
$X_0V_M(W)\asymp_{M,\Psi}xJ_0L^{-1.24}$.

The main sieve product is $\Pi_r(W)=V_M(W)(1+o(1))$ uniformly.
Also the prime harmonic estimate gives
\[
 \sum_{n\in\mathcal N_{\gamma,\gamma'}}\frac1{n-1}
       =\log\frac{\gamma'}\gamma+o(1).
\]
Here replacing $1/n$ by $1/(n-1)$ costs $o(1)$, and Mertens' theorem
evaluates the sum of $1/n$.  These facts give the main term in
\eqref{eq:conditioned-bin-count} for $r\le x^\epsilon$.  For the
omitted tail, a supported $d<2x$ has at most
$\log(2x)/(bL)=O_b(1)$ distinct prime divisors above $x^b$.
Multiplying the tail bound \eqref{eq:marked-tail} by this fixed
factor still gives $o(X_0V_M(W))$.  This proves
\eqref{eq:conditioned-bin-count}.

\subsection{Removing the least prime factor}

If the least prime factor of a supported composite $d$ belongs to
the bin $(x^\gamma,x^{\gamma'}]$, write it as $n$ and put $m=d/n$.
Then $P^-(m)\ge n>x^\gamma$, also when $n^2\mid d$.  The weight of
these composites is consequently at most
\[
 T_{\gamma,\gamma'}=
 \sum_{n\in\mathcal N_{\gamma,\gamma'}}
       \ \sum_{\substack{m\ge1\\P^-(m)>x^\gamma}}w(mn).
\]
Lemma~\ref{lem:twisted-rough-type-ii} gives the replacement
\begin{equation}\label{eq:rough-replacement}
 T_{\gamma,\gamma'}=
 \sum_{n\in\mathcal N_{\gamma,\gamma'}}
       \ \sum_{\substack{m\ge1\\P^-(m)>W\\x<mn<2x}}
       \frac{D_\gamma(\log m/L)}{LV(W)}w(mn)
       +o(X_0/L).
\end{equation}
We check the geometry, support, and total error in this application.

Partition each factor into disjoint half-open dyads
$[H_m,2H_m)$ and $[H_n,2H_n)$, restricting $n$ to its bin.  Retain
every rectangle containing a real point with $mn/x\in\operatorname{supp}\Psi$
and $x^\gamma<n\le x^{\gamma'}$.  In such a rectangle
$x/4<H_mH_n<2x$.  There are $O(L)$ rectangles: there are $O(L)$
possible dyads for $n$, and the displayed product constraint allows
only an absolute bounded number of dyads for $m$ for each of them.
On the support of $\Psi(mn/x)$, write $t=\log n/L$.  Then
\begin{equation}\label{eq:bin-exponent-range}
 \frac{\log m}{L}=1-t+O(L^{-1}),
 \qquad b\le\gamma<t\le\gamma'\le\frac12-\kappa.
\end{equation}
Passing from a point in this support to the whole $m$-dyad changes
its logarithmic exponent by at most $\log2/L$.  Therefore, for all
large $x$, the full $m$-dyad is in the exponent interval
\[
 [w_-,w_+]=\left[\frac12+\frac\kappa2,\,1-\frac b2\right].
\]
In particular $0<\gamma<w_-<w_+<1$.  Both dyad scales exceed
$x^{b/2}$ for large $x$: the $n$-scale is at least $x^b/2$, while
the $m$-scale is at least $x^{1/2+\kappa}/2$.

For the separation of the smooth weight, put
\[
 \widehat\psi(v)=\frac1{2\pi}\int_{\mathbb R}\Psi(e^s)e^{-ivs}\dd s.
\]
Fourier inversion and smooth compact support give
\begin{equation}\label{eq:reservoir-fourier-separation}
 \begin{gathered}
 \Psi(mn/x)=\int_{\mathbb R}\widehat\psi(v)x^{-iv}m^{iv}n^{iv}\dd v,\\
 \int_{\mathbb R}|\widehat\psi(v)|\dd v\ll_\Psi1,
 \qquad
 \int_{|v|>L}|\widehat\psi(v)|\dd v\ll_{A,\Psi}L^{-A}
 \end{gathered}
\end{equation}
for every fixed $A>0$.  Insert the character expansion
\eqref{eq:product-progression} to separate the progression condition.

In each retained box apply Lemma~\ref{lem:twisted-rough-type-ii}
with
\[
 \delta=b/2,\qquad C=2,\qquad D_*=6,
\]
and with the exponent interval above.  For $|v|\le L$ its first
coefficient is the difference of the two tests in
\eqref{eq:rough-replacement}, multiplied by $\chi_0(m)m^{iv}$.
The second coefficient is
$\chi_0(n)n^{iv}\one_{n\text{ prime}}$, restricted to its bin and
dyad.  It is bounded by one and is supported on $W$-rough integers,
since $n>W$.  The function $F$ is exactly the invariant function
in \eqref{eq:marked-family}.  The product condition remains in
$\Psi$, so neither coefficient has been restricted by a condition
on both factors.  The lemma applies for sufficiently large $K$
depending on these fixed parameters, independently of the band
exponents.  Its bound \eqref{eq:twisted-rough-type-ii} is
$O(xL^{-6})$ per box; integration over $|v|\le L$ costs only the
bounded $L^1$ norm in \eqref{eq:reservoir-fourier-separation}.

For $|v|>L$, the density is bounded on the fixed exponent interval
and $LV(W)\asymp L^{0.76}$.  Thus the two discrepancy coefficients
are bounded for large $x$.  The marked weight is bounded by $B_K$ by
\eqref{eq:mark-bound},
and a box has $O(H_mH_n)=O(x)$ integer pairs.  The last bound in
\eqref{eq:reservoir-fourier-separation}, with sufficiently large
fixed $A$, makes this tail $O_{K,\Psi}(xL^{-6})$ per box as well.  The
$O(L)$ boxes therefore contribute $O(xL^{-5})=o(X_0/L)$, with the
constant allowed to depend on all fixed parameters, using
\eqref{eq:marked-mass-comparability} and $J_0\ge1$.  This proves
\eqref{eq:rough-replacement}.  All occurrences of the density in
that equation are in its domain, by \eqref{eq:bin-exponent-range}.

By the joint continuity established in Section~\ref{sec:marked-family},
\eqref{eq:bin-exponent-range} implies, uniformly on the support in
\eqref{eq:rough-replacement},
\[
 D_\gamma(\log m/L)
 \le\sup_{\gamma\le t\le\gamma'}D_\gamma(1-t)+o(1).
\]
Since $n>W$, the condition $P^-(m)>W$ is equivalent to
$P^-(mn)>W$.  For a fixed divisor $n\mid d$ the factor $m=d/n$ is
unique.  We may therefore apply \eqref{eq:conditioned-bin-count}
after changing variables $d=mn$.  Together with
$V_M(W)/V(W)\to\mathfrak S_M$, this gives
\begin{equation}\label{eq:least-factor-bin-upper}
 T_{\gamma,\gamma'}\le\frac{\mathfrak S_M X_0}{L}
 \left\{
   \left(\sup_{\gamma\le t\le\gamma'}D_\gamma(1-t)\right)
                   \log\frac{\gamma'}\gamma+o(1)
 \right\}.
\end{equation}

Choose a finite partition
$b=\gamma_0<\gamma_1<\cdots<\gamma_J=1/2-\kappa$.
The set of pairs $(\gamma,1-t)$ with
$b\le\gamma\le t\le1/2-\kappa$ is a compact subset of
$\{(\gamma,w):w>\gamma>0\}$, since $1-t-\gamma\ge2\kappa$.
Joint continuity of the density is uniform there.  For every $t$ in a bin,
the bin supremum in \eqref{eq:least-factor-partition}, whose first argument
is fixed at the left endpoint, differs from $D_t(1-t)$ by a quantity
tending uniformly to zero with the mesh.  The measure $\dd t/t$ has finite total mass
$\log((1/2-\kappa)/b)$.  We may consequently choose a sufficiently
fine fixed partition for which
\begin{equation}\label{eq:least-factor-partition}
 \begin{split}
 \sum_{j=1}^J
  \left(\sup_{\gamma_{j-1}\le t\le\gamma_j}
                     D_{\gamma_{j-1}}(1-t)\right)
             \log\frac{\gamma_j}{\gamma_{j-1}}
 &\le\int_b^{1/2-\kappa}D_t(1-t)\frac{\dd t}{t}+\frac18\\
 &\le D_b(1)-1+\frac18.
 \end{split}
\end{equation}
The last inequality uses nonnegativity of $D_t$ and
Input~\ref{input:density}.

Put
\[
 U_\kappa=\sum_{\substack{d\ge2\\P^-(d)>x^{1/2-\kappa}}}w(d).
\]
Every supported integer counted in \eqref{eq:initial-rough-mass}
but not in the sum defining $U_\kappa$ is composite: it exceeds $x$,
whereas its
least prime factor is at most $x^{1/2-\kappa}$.  That factor lies
in one of the bins of our partition.  Subtracting
\eqref{eq:least-factor-bin-upper} for these finitely many bins from
\eqref{eq:initial-rough-mass}, and using
\eqref{eq:least-factor-partition} and
\eqref{eq:density-integral-bound}, yields
\begin{equation}\label{eq:post-bin-mass}
 \frac{LU_\kappa}{\mathfrak S_M X_0}
 \ge1-\frac18-\frac{e^{-\gamma_E}}{b}
       \left(C_se^{-c_s/b}+C_de^{-c_d/b}\right)+o(1).
\end{equation}
Here the lower bound uses
$\frac{e^{-\gamma_E}}b(1-C_se^{-c_s/b})-
 (D_b(1)-1+1/8)$; the stated upper bound on $D_b(1)$ gives precisely
the two exponential losses in \eqref{eq:post-bin-mass}.

\subsection{Balanced composites and the parameter choices}

It remains to bound the composite part of $U_\kappa$.  The final choice of
$\kappa$ precedes those of $b$ and $K$, so the coefficient in this bound
must be independent of all three.  We achieve this by sieving pairs for
each fixed $r$ before summing the marks.  Such a composite has exactly two
prime factors counted with multiplicity.
Indeed three factors, each greater than $x^{1/2-\kappa}$, would have
product at least $x^{3/2-3\kappa}>2x$ for large $x$.  Write the two
factors as $m,n$.  Each is greater than $x^{1/2-\kappa}$ and less
than $2x^{1/2+\kappa}$, so for large $x$ both belong to
\[
 [x^{1/2-2\kappa},x^{1/2+2\kappa}].
\]
Use ordered pairs $(m,n)$ to bound their weight; this counts a square
once and may count other semiprimes twice, which is harmless for an
upper bound.

Cover these pairs by disjoint dyadic boxes
$[M_1,2M_1)\times[M_2,2M_2)$ meeting the displayed factor range and
$x<mn<2x$.  Every retained box satisfies
\begin{equation}\label{eq:balanced-box-range}
 \frac x4<M_1M_2<2x,
 \qquad M_i\le2x^{0.54}\quad(i=1,2).
\end{equation}
The latter is a uniform relaxed bound because $\kappa<0.01$.
There are at most $C_{\mathrm{dyad}}(\kappa L+1)$ boxes, for an
absolute constant $C_{\mathrm{dyad}}$.  To see this, the logarithmic
length of the factor range is $4\kappa L$, so it meets
$O(\kappa L+1)$ dyads for $m$; for each of them the first inequality
in \eqref{eq:balanced-box-range} permits only an absolute bounded
number of dyads for $n$.

Fix one box and a group integer $r$ in the full support, so
$r\le x^{0.11}$.  Put $k_r=4r$ and $z=x^{0.002}$.  A prime pair
contributing to the weight through this $r$ obeys
\begin{equation}\label{eq:balanced-necessary-conditions}
 mn\equiv1\pmod{k_r},
 \qquad
 mn\not\equiv0,1\pmod j
 \quad(j\le z\text{ an odd prime},\ j\nmid r).
\end{equation}
In fact $m,n>x^{1/2-2\kappa}>z$ and
$Q=(mn-1)/(4r)>x^{0.9}>z$ are primes.  Thus $mn\equiv0\pmod j$
would force $j$ to be one of the prime factors $m,n$, and, since
$j\nmid4r$, $mn\equiv1\pmod j$ would force $j=Q$.  Both are
impossible.  We discard the progression condition modulo $M$ to
obtain an upper bound.

Use as a finite unweighted family all integer pairs in the box
satisfying the first congruence in
\eqref{eq:balanced-necessary-conditions}, and designate its odd
primes $j\le z$ not dividing $r$, with bad condition
$mn\equiv0\text{ or }1\pmod j$.  All sums and products indexed by
$j$ in this balanced sieve are over primes.  For a squarefree product
$\ell$ of these primes set
\[
 \rho(\ell)=\prod_{j\mid\ell}(3j-2),\qquad \rho(1)=1.
\]
There are $\varphi(k_r)$ residue pairs modulo $k_r$ with product
one.  At $j$ there are $2j-1$ residue pairs with product zero and
$j-1$ with product one, and these two sets are disjoint.  Since
$(k_r,\ell)=1$, the Chinese remainder theorem gives exactly
$\varphi(k_r)\rho(\ell)$ residue pairs modulo $k_r\ell$ satisfying
the base congruence and all the designated bad conditions at
primes dividing $\ell$.

Each residue pair modulo $k_r\ell$ occurs in the box
\[
 \frac{M_1M_2}{(k_r\ell)^2}
     +O\left(\frac{M_1+M_2}{k_r\ell}+1\right)
\]
times.  Thus its intersection count has the form $Y_rg(\ell)+R(\ell)$,
where
\begin{equation}\label{eq:balanced-sieve-data}
 Y_r=\frac{M_1M_2\varphi(k_r)}{k_r^2},
 \qquad g(\ell)=\frac{\rho(\ell)}{\ell^2}.
\end{equation}
\begin{equation}\label{eq:balanced-lattice-error}
 \begin{split}
 |R(\ell)|
 &\ll\varphi(k_r)\rho(\ell)
       \left(\frac{M_1+M_2}{k_r\ell}+1\right)\\
 &\ll (M_1+M_2)\ell+k_r\ell^2.
 \end{split}
\end{equation}
The last inequality uses $\varphi(k_r)\le k_r$ and
$\rho(\ell)\le\ell^2$, because $3j-2\le j^2$ for every odd prime
$j$.  This count also includes $\ell=1$, so its boundary error for
the base family is included in $R(1)$.

For a designated prime,
$g(j)=(3j-2)/j^2\le7/9$ and $g(j)\le3/j$.  Thus the local gap and
block sums in Input~\ref{input:block} again have absolute constants.
Its $H=2$ upper bound uses only $\ell\le z^{10}=x^{0.02}$.
Summing \eqref{eq:balanced-lattice-error} over this range, bounding
the designated squarefree products by all positive integers, gives
\begin{align}
 \sum_{\ell\le x^{0.02}}|R(\ell)|
 &\ll (M_1+M_2)\sum_{\ell\le x^{0.02}}\ell
        +k_r\sum_{\ell\le x^{0.02}}\ell^2\notag\\
 &\ll x^{0.54+0.04}+r x^{0.06}
 \ll x^{0.58}+x^{0.17}.\label{eq:balanced-remainder-sum}
\end{align}
The notation in the left-hand sum restricts to the designated
products for which $R$ is defined.  The exponent margins compared
with $x/r$ are uniform throughout $r\le x^{0.11}$:
\[
 \frac{x^{0.58}}{x/r}=r x^{-0.42}\le x^{-0.31},
 \qquad
 \frac{r x^{0.06}}{x/r}=r^2x^{-0.94}\le x^{-0.72}.
\]
Both fixed power savings absorb $L^4$.  It follows that
\eqref{eq:balanced-remainder-sum} is $o((x/r)L^{-4})$ uniformly in
every supported $r$, with no dependence on $K$ in this assertion.

For the main product the exact local identity is
\[
 1-g(j)=\left(1-\frac1j\right)^3
            \left(1-\frac1{(j-1)^2}\right).
\]
The last factor is at most one.  The missing factor at two costs
only the factor $8$ when compared with $V(z)^3$.  The missing
factors at primes dividing $r$ cost at most
\[
 \prod_{j\mid r}\left(1-\frac1j\right)^{-3}
  =\exp\left(O\left(\sum_{j\mid r}\frac1j\right)\right)=1+o(1)
\]
uniformly, since on the full support
$\sum_{j\mid r}j^{-1}\le0.11L^{0.9}e^{-L^{0.1}}$ by
\eqref{eq:group-reciprocal-divisors}.  Consequently
\begin{equation}\label{eq:balanced-sieve-product}
 \prod_{\substack{j\le z\\j\ \mathrm{odd\ prime}\\j\nmid r}}
        (1-g(j))\ll V(z)^3\ll L^{-3},
\end{equation}
with an absolute constant.  The numerical exponent $0.002$ is
fixed, so the last estimate follows from Mertens' theorem with a
fixed constant.

By \eqref{eq:balanced-box-range}, $Y_r\ll x/r$.  Combining
\eqref{eq:block-sieve-upper}, \eqref{eq:balanced-remainder-sum},
and \eqref{eq:balanced-sieve-product} bounds the number of prime
pairs in this box contributing through $r$ by
\[
 O\bigl((x/r)L^{-3}\bigr),
\]
uniformly for all supported $r$.  Its constant is absolute; the
power-saving error above has been absorbed for sufficiently large
$x$.  The smooth weight is at most one.  Multiplying by
$\mathcal W(r)$ and summing over the full support of $r$, without
discarding a tail, therefore bounds the weighted prime pairs in
one box by $C_{\mathrm{pair}}xJ_0/L^3$ with an absolute
$C_{\mathrm{pair}}$.  The unique
representation of the support justifies this assignment of each
weighted pair to its group integer.

Let $\mathcal C_\kappa$ denote the weight of the composites counted by
$U_\kappa$.  Summing the last estimate over the boxes gives
\begin{equation}\label{eq:balanced-composite-mass}
 \mathcal C_\kappa
 \le C_2(\kappa+L^{-1})\frac{X_0}{L}.
\end{equation}
Here $C_2$ depends only on $M,\Psi$ and the fixed dyadic comparison
conventions, and is independent of $\kappa,b,K$.  Indeed the
preceding bound before converting the normalization is
$C_{\mathrm{pair}}C_{\mathrm{dyad}}(\kappa L+1)xJ_0/L^3$.
The lower comparison in \eqref{eq:marked-mass-comparability} is
$X_0\ge c_{M,\Psi}xJ_0/L$ for large $x$, with
$c_{M,\Psi}>0$ independent of $K$.  We may take
$C_2=C_{\mathrm{pair}}C_{\mathrm{dyad}}/c_{M,\Psi}$.  The threshold for $x$ may
depend on all the fixed parameters; this does not change the stated
independence of the constant.

Every supported prime is counted by $U_\kappa$.  Subtracting its
composite part using \eqref{eq:post-bin-mass} and
\eqref{eq:balanced-composite-mass} yields
\begin{equation}\label{eq:final-prime-mass}
 \begin{split}
 \frac{L}{\mathfrak S_M X_0}\sum_{d\text{ prime}}w(d)
 &\ge1-\frac18-\frac{e^{-\gamma_E}}{b}
       \left(C_se^{-c_s/b}+C_de^{-c_d/b}\right)\\
 &\quad-\frac{C_2}{\mathfrak S_M}(\kappa+L^{-1})+o(1).
 \end{split}
\end{equation}
We now choose the parameters.  With $M$ fixed, the function $\Psi$ and
the dyadic convention have already fixed $C_2$, while
$\mathfrak S_M>0$ depends only on $M$.  Choose
$0<\kappa<0.01$ so small that
$C_2\kappa/\mathfrak S_M<1/8$, and then choose
$0<\epsilon<\kappa$.  Next choose fixed $b>0$ sufficiently small
for Inputs~\ref{input:block} and~\ref{input:density}, for
\eqref{eq:initial-sieve-level}, and so that
\[
 \frac{e^{-\gamma_E}}{b}
       \left(C_se^{-c_s/b}+C_de^{-c_d/b}\right)<\frac18.
\]
This is possible because $b^{-1}e^{-c/b}\to0$ as $b\downarrow0$
for each $c>0$.  Fix a finite partition satisfying
\eqref{eq:least-factor-partition}.  For its finitely many left
endpoints $\gamma$, use the fixed Type II choices
$\delta=b/2$, $C=2$, $D_*=6$ and the exponent interval already
specified.  Now choose $K$ sufficiently large for
Lemma~\ref{lem:twisted-rough-type-ii} in every one of these
applications.  The lower bound on $K$ is independent of the band
exponents, so choose any fixed
$0.1<a_1<\cdots<a_K<0.2$ afterwards.  Finally let $x$ exceed all
thresholds for these fixed choices.

The three fixed losses in \eqref{eq:final-prime-mass} are then at
most $1/8$ each, and the term $C_2/(\mathfrak S_M L)$ and the
$o(1)$ term tend to zero.  In particular, for all sufficiently
large $x$,
\[
 \sum_{d\text{ prime}}w(d)
       \ge\frac{\mathfrak S_M X_0}{2L}
       \gg_{M,\Psi}\frac{xJ_0}{L^2}
       \gg_{M,\Psi}\frac{x}{L^2}.
\]
By \eqref{eq:marked-family-support}, each prime with positive weight
has the required support and progression properties, and its
weight is at most $B_K$.  Dividing the prime mass by this fixed
bound gives the required lower bound for distinct primes.
This proves Proposition~\ref{prop:reservoir}.
